%% LyX 2.4.4 created this file.  For more info, see https://www.lyx.org/.
%% Do not edit unless you really know what you are doing.
\documentclass[english]{article}
\usepackage[T1]{fontenc}
\usepackage[latin9]{inputenc}
\usepackage{cancel}
\usepackage{geometry}
\geometry{verbose,tmargin=1in,bmargin=1in,lmargin=1in,rmargin=1in}
\PassOptionsToPackage{normalem}{ulem}
\usepackage{ulem}

\makeatletter

%%%%%%%%%%%%%%%%%%%%%%%%%%%%%% LyX specific LaTeX commands.
%% Because html converters don't know tabularnewline
\providecommand{\tabularnewline}{\\}

%%%%%%%%%%%%%%%%%%%%%%%%%%%%%% User specified LaTeX commands.
\usepackage{hyperref}
\usepackage[usenames,dvipsnames]{xcolor} %adds additional color options
\hypersetup{colorlinks=true, urlcolor=Blue}%Blue

\usepackage{sectsty} %%one way to change section header font sizes.
\sectionfont{\large}

\usepackage{titlesec}
\titlespacing*{\section}{0pt}{10pt}{0pt}

\usepackage{multicol}

\usepackage{tikz}
\usetikzlibrary{plotmarks}
\usepackage{pgfplots}
\pgfplotsset{compat=newest}

\usepackage{calc}
\usepackage[nomessages]{fp}

\usepackage{advdate}    % Advancing/saving dates
\usepackage{datetime}   % Dates formatting
\usepackage{datenumber} % Counters for dates
%\longdate %\usdate %\DM %\MY
\newdateformat{MD}{\monthname[\THEMONTH] \ordinal{DAY}}
% Added by lyx2lyx
\setlength{\parskip}{\medskipamount}
\setlength{\parindent}{0pt}

\makeatother

\usepackage{babel}
\begin{document}
\begin{center}
{\Large\textbf{Johnson C. Smith University}}\textbf{}\\
\textbf{Introductory Abstract Algebra I~~|~~MTH 335 A~\,|~~Spring
2013}\\
\textbf{\rule{1\columnwidth}{1pt}}
\par\end{center}

\vspace{-.6cm}

\begin{center}
\emph{Pure mathematics is, in its way, the poetry of logical ideas}.
\\
\textasciitilde Albert Einstein
\par\end{center}

\begin{center}
\emph{A problem worthy of attack, proves its worth by fighting back!}
\\
\textasciitilde Paul Erdos
\par\end{center}

\section*{Instructor Information}

\textbf{Name: }Dr.\,James Ashe\\
\textbf{Office location: }307 Davis Hall (back); office hours will
also be held in the Math Lab, EDU 115\\
\textbf{Email: }\href{mailto:jashe@jcsu.edu}{jashe@jcsu.edu}\\
\textbf{Office phone: }(704)\,378-1037\\
\textbf{Office hours: }MWF 10\textendash 11am and 12-1pm; By appointment
only, MW 8-9am and 2-3pm.

\section*{Course Description}

This is a \textbf{3 credit hour}\textbf{\emph{ }}course; the first
in a two-course sequence which are offered for students from the sciences,
social sciences and engineering. Topics covered: Elementary set theory
and logic. Mappings, groups, rings, integral domains, fields, and
polynomials. Prerequisite: Mathematics 231 or consent of Department.

Instructional methods will consist of lectures, in-class discussions,
exams, and student presentations. Other methods may be adopted according
to the particular needs of the class.

\section*{Course Objectives}

The main objective of this course is to develop mathematical reasoning
by learning how to read, write, and present proofs. Students completing
this course will demonstrate knowledge of the fundamental concepts
in the core areas listed in the above section. Additionally, students
should demonstrate the ability for abstract reasoning, rigor, and
logical thinking.

\section*{Course Materials }

\textbf{Suggested Textbook: }There is no required textbook. Below
is a list of recommended reference books. The first three are free.\textbf{}\\
\emph{\href{http://abstract.ups.edu/download/aata-20120811.pdf}{\emph{Abstract Algebra: Theory and Applications} by Thomas W.\,Judson}}\\
\emph{\href{http://www.csun.edu/~asethura/giaa/giaa.pdf}{\emph{A Gentle Introduction to Abstract Algebra} B.\,A\,Sethuraman}}\\
\emph{\href{http://www.math.uiuc.edu/~r-ash/Algebra.html}{\emph{Abstract Algebra: The Basic Graduate Year} by Robert Ashe}}\\
\emph{\href{http://www.amazon.com/Contemporary-Abstract-Algebra-Fifth-Edition/dp/0618122141/ref=sr_1_sc_1?ie=UTF8&qid=1359086349&sr=8-1-spell&keywords=algerba+gallian}{\emph{Contemporary Abstract Algebra} by Joseph A.\,Gallian}}

\textbf{Calculator: }Feel free to bring one if you want.

\section*{Grades}

\textbf{Grading scale: $\mbox{A}\geq90>\mbox{B}\geq80>\mbox{C}\geq70>\mbox{D}\geq60>\mbox{F}$}

\textbf{Evaluation: }Grades will be taken from homework, quizzes,
and Exams as follows:\textbf{}\\
Homework$=50\%$; Presentations$=15\%$; Midterm$=15\%$; Final $=20\%$

\textbf{Homework}: Regular homework assignments will be given. Students
can collaborate on homework assignments, however each student should
hand in their own assignments.\textbf{ }The lowest homework score
will be dropped.\textbf{}\\
\textbf{Presentations: }Each student will be given one problem to
present to the class from each homework assignment. The date of presentation
will typically precede the homework due date. You will be assessed
on both your ability to work the problem and explain it.\textbf{}\\
\textbf{Midterm and Final Exam: }Both the midterm and final will be
cumulative exams. The midterm will be given on Friday, March 22, and
the final will be given in our normal class room on Friday, May 17
from 1-3pm; see the exam schedule:\\
\href{http://www.jcsu.edu/happenings/exam-schedule/spring-exam}{www.jcsu.edu/happenings/exam-schedule/spring-exam}.\emph{
If beneficial, the final exam score will replace the midterm exam
score.}

\section*{Class Policies}

JCSU policies regarding attendance, academic integrity, grades, and
change of enrollment will be enforced. It is the student\textquoteright s
responsibility to be aware of these policies and of related deadlines
as well as any announcements or syllabus adjustments (written or oral)
made in class. Please see the student handbook for details, \href{http://www.jcsu.edu/docs/handbook.pdf}{jcsu.edu/docs/handbook.pdf}.

\textbf{Attendance: }Class attendance is required for all JCSU students.
Each student is allowed as many hours of absence per term as credit
hours(s) received (not to exceed 4) for the class. Students who exceed
the maximum number of absences may receive a failing grade for the
course. 

Students who may miss classes while representing the University in
an official capacity are exempt from regulations governing absences.
However, absence from class for official University business does
not relieve the student from responsibility for any class assignments
that may be missed during the period of absence. In such cases, \emph{it
is the student's responsibility to make arrangements to make up any
work that will be missed }\emph{\uline{before the assignment/quiz/exam
is missed.}}

\textbf{Late Homework: }An assignment will have 25\% deducted for
each day it is late.\textbf{}\\
\textbf{Missed presentations:} A student will be allowed one make-up
for a missed presentation. Presentations can be rescheduled If a University
sponsored event conflicts with the presentation date, and advance
notice of the absence was given by the student. \\
\textbf{Missed exams: }Missed exams may be made up at the discretion
of the Instructor. \uline{(1) you must notify me} \uline{by email
or phone within} 24 hours of the missed exam AND \uline{(2) you must
have no more than 3 unexcused absences in the class}. If BOTH of these
conditions are met, a make-up will be given. However unless a \uline{documented
}excuse is provided, 10\% will automatically be deducted from the
make-up. If you will be gone for any reason, it will be \emph{your
}responsibility to inform me at least a week in advance to schedule
to take the exam early. 

\section*{Other}

\noindent\begin{minipage}[t]{1\columnwidth}%
\begin{tabular}{rcl}
\textbf{Important Dates} &  & \tabularnewline
\formatdate{22}{3}{2013} &  & Exam 4\tabularnewline
\formatdate{17}{5}{2013} &  & Final Exam\tabularnewline
 &  & \tabularnewline
\end{tabular}%
\end{minipage}

\textbf{Student support services: }If you need course adaptations
or accommodations for a documented disability please contact the Office
of Disability Services. \\
\href{http://www.jcsu.edu/directory/federal-requirements/general-institutional-information/jcsu-disability-services}{jcsu.edu/directory/federal-requirements/general-institutional-information/jcsu-disability-services},
\\
phone: (704)\,398-1116.

The Office of Counseling Services provides confidential intakes and
assessments, personal counseling and referrals to appropriate resources
to all enrolled students on a walk-in or appointment basis. \\
\href{http://www.jcsu.edu/admissions/future_students/meet_your_counselor}{jcsu.edu/admissions/future\_students/meet\_your\_counselor},
phone (704)\,378-1044.

\textbf{Disclaimer: }While I do not foresee any significant changes
to the items in this syllabus during the semester, I reserve the right
to make appropriate changes. Students will be notified of any such
changes in writing.

JCSU's Fall 2013 academic calendar can be found here: \href{http://www.jcsu.edu/happenings/academic-calendar}{jcsu.edu/happenings/academic-calendar}

\vfill

\begin{center}
\emph{The mathematician does not study pure mathematics because it
is useful; he studies it because he delights in it and he delights
in it because it is beautiful.}\\
\textasciitilde Poincare
\par\end{center}

\vfill
\end{document}
